\chapter{Overview}

\begin{center}
	\includegraphics[width=0.5\textwidth]{Images/aberfitness.png}
\end{center}

\par
\textit{Aber Fitness} is a web application developed using Microsoft's \textit{.NET Core} and Oracle's \textit{Java Enterprise Edition} (henceforth referred to as Java EE). The project aims to provide a service to encourage fitness and promote engagement with sporting activities amongst the users of the application, offering functionality such as:

\begin{itemize}
\item Graphing fitness data gathered by owners of \textit{Fitbit} devices, or through manual input.
\item The ability to take part in challenges based on that fitness data.
\item Providing a Sport ladder system, which has tight integration with a bespoke facility booking system.
\end{itemize}

\textit{Aber Fitness} aims to offer everything that would be needed by a sporty and active person in order to bring their sporting activities into a digital platform and also to enhance their use of devices they already own, such as \textit{Fitbit} devices or smart watches such as the \textit{Apple Watch}.

\par
At launch, the system will ingest activity data automatically from \textit{Fitbit}, and be open to the addition of other health data provider services at a later date due to the modular nature of the data ingest system. Once normalised this activity data will be used throughout the various subsystems of \textit{Aber Fitness}, providing users with functionality such as a dashboard overview of their activity over approximately a month, as well as integrating tightly with the challenges system to add a competitive aspect to the system, with the aim to keep users engaged with both the platform itself and keeping fit in general. 

\par
Aber Fitness will also feature a number of supporting tools, such as:

\begin{itemize}
\item A user management system with two-factor authentication.
\item The ability for users to view a comprehensive audit trail of when and how their data was used and accessed.
\item A global layout template with a dynamic sidebar, which can be adjusted by administrators and those changes are instantly propagated to other microservices.
\end{itemize}

\par
The platform is also designed to be scalable and easily deployed through the use of \textit{Docker}, as each aspect of the system has been segregated into individual microservices. 